\documentclass[11pt]{article}
\usepackage[margin=1in]{geometry}
\usepackage{amsmath,amssymb,amsthm,mathtools,microtype}
\usepackage[hidelinks]{hyperref}
\newtheorem{theorem}{Theorem}
\newtheorem{corollary}[theorem]{Corollary}
\newtheorem{remark}[theorem]{Remark}
\newcommand{\R}{\mathbb R}
\newcommand{\cC}{C}
\newcommand{\HH}{\mathbb H}
\title{A smooth-conical partial answer for smooth singular cochains on real algebraic sets}
\author{Autonomous proof attempt for arXiv:2208.11431}
\date{11 August 2026}

\begin{document}
\maketitle

\begin{abstract}
The source asks whether the coaugmentation from the constant sheaf to the
sheafification of smooth singular cochains on the real spectrum of every
finitely generated real algebra is a quasi-isomorphism.  We prove the desired
conclusion whenever the underlying predifferentiable space is locally smoothly
contractible.  This gives a concrete positive answer for locally
smooth-conical algebraic sets, in particular for isolated positively
weighted-homogeneous singularities and for normal-crossing algebraic sets.
The argument is stalkwise and uses the smooth singular prism operator.  We
also isolate why the known Lipschitz contraction does not settle the general
question.
\end{abstract}

\section{The exact question}

Let $B$ be a finitely generated real algebra and let $X=\operatorname{spec}B$
carry the predifferentiable structure induced by a distinguished algebraic
embedding into Euclidean space.  Write $\cC^q_{\mathrm{sm}}(U)$ for the dual
of the real vector space of smooth singular $q$-chains in $U$, and let
$${}^{+}\cC^\bullet_{\mathrm{sm},X}$$
be the sheafification of this presheaf complex.  The source asks, immediately
after its Diagram~(7), whether
\begin{equation}\label{eq:epsilon}
 \epsilon:\underline{\R}_X[0]\longrightarrow
 {}^{+}\cC^\bullet_{\mathrm{sm},X}
\end{equation}
is a quasi-isomorphism, equivalently whether the induced map on
hypercohomology is an isomorphism.  The source proves the Lipschitz analogue
using local Lipschitz contractibility of algebraic sets, but does not know the
smooth analogue.

\section{A local criterion}

Call a predifferentiable space $X$ \emph{locally smoothly contractible} if,
for every $x\in X$ and every neighborhood $U$ of $x$, there are a neighborhood
$V\subset U$ of $x$ and a smooth map
$$H:V\times[0,1]\longrightarrow U$$
such that $H(v,0)=v$ and $H(v,1)=x$.

\begin{theorem}\label{thm:criterion}
If $X$ is locally smoothly contractible, then the coaugmentation
\eqref{eq:epsilon} is a quasi-isomorphism.  Consequently it induces
isomorphisms
$$
 \HH^k(X,\underline{\R}_X[0])\ \cong\
 \HH^k(X,{}^{+}\cC^\bullet_{\mathrm{sm},X})
 \qquad(k\geq0).
$$
\end{theorem}

\begin{proof}
The stalk of a sheafification equals the stalk of the original presheaf.
Moreover, filtered colimits of real vector spaces are exact.  We may therefore
compute the cohomology sheaves of
$ {}^{+}\cC^\bullet_{\mathrm{sm},X}$ by taking the cohomology of its stalk
complexes.

Fix $x\in X$.  A germ in the stalk is represented by a smooth singular
cochain on some neighborhood $U$ of $x$.  After shrinking $U$ if necessary,
a cocycle germ is represented by an actual cocycle.  Choose $V$ and $H$ as in
the definition.  The standard prism decomposition of
$\Delta^q\times[0,1]$ is affine on finitely many simplices.  Composing its
pieces with $H\circ(\sigma\times\mathrm{id})$ therefore defines the usual
smooth singular chain homotopy between the inclusion $V\hookrightarrow U$
and the constant map $V\to\{x\}\hookrightarrow U$.  Thus restriction from
$U$ to $V$ factors on smooth singular cohomology through the cohomology of a
point.  It is zero in positive degrees.  In degree zero the same argument,
or smooth path connectedness of $V$ supplied by $H$, identifies the kernel of
the coboundary with the constant functions.

It follows that every positive-degree cohomology germ vanishes and that the
degree-zero cohomology sheaf is $\underline{\R}_X$.  Hence \eqref{eq:epsilon}
is a quasi-isomorphism.  Hypercohomology preserves quasi-isomorphisms.
\end{proof}

\section{Weighted cones}

For positive integers $w_1,\dots,w_r$, write
$$
 \delta_t(z_1,\dots,z_r)
   =(t^{w_1}z_1,\dots,t^{w_r}z_r),\qquad 0\leq t\leq1.
$$
A subset $C\subset\R^r$ is a positive weighted cone if
$\delta_t(C)\subset C$ for all $t\in[0,1]$.  The map $(z,t)\mapsto\delta_tz$
is polynomial and hence smooth, including at $t=0$.

Say that $X$ is \emph{locally smooth-conical} if every point has a
predifferentiable neighborhood smoothly isomorphic to a relatively open
neighborhood of $(0,0)$ in $M\times C$, where $M$ is a smooth manifold chart
centered at $0$ and $C$ is a positive weighted cone.

\begin{corollary}\label{cor:cone}
For every locally smooth-conical real algebraic set $X$, the coaugmentation
\eqref{eq:epsilon} is a quasi-isomorphism and the questioned
hypercohomology map is an isomorphism in every degree.
\end{corollary}

\begin{proof}
Shrink the manifold chart to a ball and contract it linearly to its center.
Simultaneously contract $C$ by $\delta_t$.  On a sufficiently small product
neighborhood this is a smooth contraction.  Transfer it through the local
smooth isomorphism and apply Theorem~\ref{thm:criterion}.
\end{proof}

\begin{corollary}\label{cor:examples}
The conclusion holds in each of the following cases.
\begin{enumerate}
\item $X$ has isolated singularities and, at every singular point, its germ is
smoothly equivalent to a positive weighted cone.  In particular this covers
isolated positively weighted-homogeneous real algebraic singularities.
\item $X$ has normal-crossing local models (more generally, local models that
are products of a manifold chart with a union of coordinate subspaces).
\end{enumerate}
\end{corollary}

\begin{proof}
Away from the stated singularities, $X$ is a smooth manifold and is locally
smoothly contractible.  At a weighted-conical singularity use
Corollary~\ref{cor:cone}.  A union of coordinate subspaces is invariant under
ordinary radial dilation, so normal-crossing charts are the weight-one case.
\end{proof}

For example, the cusp $y^2=x^3$ contracts by
$(x,y)\mapsto(t^2x,t^3y)$, while a node or a union of coordinate subspaces
contracts by ordinary radial scaling.  These singular examples therefore do
not create extra local smooth singular cohomology.

\section{Why this does not settle the general case}

The source's local contraction is Lipschitz and semialgebraic.  Its
composition with an arbitrary ambient-smooth simplex need not be smooth:
the simplex can meet the contraction's nonsmooth strata along a complicated
closed set.  A flattened semialgebraic triangulation can make prescribed
semialgebraic simplices smooth at their faces, but injectivity would require a
relative smoothing theorem fixing an arbitrary smooth boundary cycle.
Neither a $C^1$ semialgebraic triangulation nor semialgebraic-current homology
supplies that step.

Likewise, a Whitney algebraic set has a conically smooth stratified structure,
but a conically smooth contraction need not be smooth for the ambient-subspace
predifferentiable structure used in the question.  Thus the argument above is
a genuine partial result: it proves the desired map is an isomorphism on a
substantial explicit class while leaving the arbitrary algebraic germ open.

\section*{Source and verification note}

The exact open sentence occurs after Diagram~(7) in Igor Baskov,
\emph{The de Rham cohomology of soft function algebras}, arXiv:2208.11431.
The literature sweep also checked
arXiv:1505.03970 (semialgebraic $C^1$ triangulations), arXiv:1404.4796
(semialgebraic current homology), and arXiv:2202.00131 (diffeological smooth
singular complexes); none supplies the missing relative ambient-smooth
comparison.

\end{document}
